% Appendix: Theorem 1.4 with explicit constants, the periodization lemmas (T2, G4.2) and Lemma EX.

\section{Companion results and analytic estimates}\label{sec:dd}

\subsection{Companion results for the reversed weight order}\label{sec:swapped}
The definition of the amplitude norm in \cite{S} and the proofs there take the two weights in different orders, and
\cite[Proposition~7.1]{S} states that the two orders give equivalent norms. For
$\beta>\tfrac12$ they do not: the two orders define different classes (Remark~\ref{pc:classes}). The principal result for the printed class is Theorem~\ref{pc:main}. In this subsection
\[
\norm a:=\norm{\ue^{\alpha\Lambda}\ue^{\beta\langle d\rangle}a}_{L^2(\Hb)} ,
\]
the order used in the proofs of \cite{S}. We call the amplitudes with $\norm a<\infty$ the \emph{swapped class}, and those with
$\norm a_{\alpha,\beta}<\infty$ the \emph{printed class}.
By Lemma~\ref{lem:m2w} below, a finite norm $\norm a$ at radius $\alpha$ gives a finite norm with the weight
$\ue^{\beta\tilde\psi}$ in place of $\ue^{\beta\langle d\rangle}$, at the radius $c_0\alpha$ for a fixed
$c_0=c_0(\beta)\in(0,1]$; the proofs use $\tilde\psi$.
\begin{theoremR}\label{thm:main}
Let $\Gamma$ be any Fuchsian group, $\beta>1$ and $\alpha\in(0,\alpha_0]$ constants ($\alpha_0<\pi/4$ fixed).
There are $C,\delta>0$ such that for all $0<\hbar\le1$, every $b\in\partial\Hb$ and every $a$ with $\norm a<\infty$,
\[
\norm{u^{(0)}(t)_\Gamma-\cU(t)[a\ue^{i\varphi_b/\hbar}]_\Gamma}_{L^2(\Hb/\Gamma)}\le C\hbar t\norm a,\qquad 0\le t\le\delta\hbar^{-1/2}.
\]
\end{theoremR}
For $\alpha>\alpha_0$ the class with radius $\alpha$ is contained in the class with radius $\alpha_0$
(in the swapped order this nesting is trivial), so the theorem holds for every constant $\alpha>0$,
as in \cite{S}. By Proposition~\ref{prop:g5} the order $\hbar t$ cannot be improved.

\begin{proof}
This is the case $K=N=0$ of Theorem~\ref{thm:dd14} below, with $\kappa:=\alpha^3$ and with $\alpha_0$ as here: then
$\alpha^3\ge\kappa\hbar$ for all $0<\hbar\le1$, $u^{(0)}$ is the approximation of order $0$, and the bound of Theorem~\ref{thm:dd14} reads
$C\alpha^{-2}\hbar t\norm a$ for $0\le t\le\delta\alpha^{1/2}\hbar^{-1/2}$. The constants of the present theorem are
$C\alpha^{-2}$ and $\delta\alpha^{1/2}$, with $C$ and $\delta$ the constants of Theorem~\ref{thm:dd14} for $\kappa=\alpha^3$; they depend on the fixed $\alpha$.
\end{proof}

\begin{remark}
No Volterra resummation and no Sobolev inequality on $M$ is used, so the proof holds for every Fuchsian
group; the cusp issue of Remark~\ref{rem:g4} does not arise. Since Theorem~\ref{thm:dd14} holds for every $\beta>1/2$, so does
the present theorem; we state it with $\beta>1$ as in \cite{S}.
\end{remark}

\medskip\noindent\textbf{The estimate of order $K$ and its corollaries.}
Let $u^{(K)}$ be the approximation of order $K$ of Section~\ref{sec:setup}.
\begin{theoremR}\label{thm:dd14}
Fix $\beta>1/2$, $\kappa>0$ and $\alpha_0\in(0,\min(1,\pi/4))$. There is $\delta>0$, and for every $P<\infty$ there is $C>0$,
such that for every Fuchsian group $\Gamma$ with $\sup_z\sum_{\gamma\in\Gamma}\ue^{-(\beta+1/2)\,d(x_0,\gamma z)}\le P$, all $0<\hbar\le1$, all $\alpha\in(0,\alpha_0]$ with $\alpha^3\ge\kappa\hbar$, all $K,N\in\mathbb N_0$, all $b\in\partial\Hb$, all $a$
with $\norm a<\infty$, and all $0\le t\le\delta(\alpha/\hbar)^{1/2}$,
\[
\norm{\Delta^N\big[\cU(t)[a\ue^{i\varphi_b/\hbar}]_\Gamma-u^{(K)}(t)_\Gamma\big]}_{L^2(\Hb/\Gamma)}
\le C^{N+K+1}N^{2N}K!\,(\alpha\hbar)^{-2N}\Big(\frac{\hbar t}{\alpha^2}\Big)^{K+1}\norm a ,
\]
with $N^{2N}:=1$ at $N=0$.
(The supremum is finite for every Fuchsian group, by Lemma~\ref{lem:p1} at the exponent $\beta+\tfrac12>1$. So $\delta$ does not
depend on $\Gamma$, and $C$ depends on $\Gamma$ only through this Poincar\'e series bound; see the end of the proof of
Theorem~\ref{thm:dd-appx}.)
\end{theoremR}
On Schubert's window $t\le\delta\alpha^{3/2}\hbar^{-1/2}$ this implies the bound \cite[(1.30)]{S} for the norm $\norm a$,
with an extra factor $(\alpha\hbar)^4\hbar t/\alpha^2\le1$. The proof is in Section~\ref{subsec:dd-thm}.

\begin{corollaryR}\label{cor:dd15}
Under the hypotheses of Theorem~\ref{thm:dd14}, with $\delta\le\kappa^{1/2}/\max(1,C)$, put $\delta_K=C\hbar t/\alpha^2$ and let $\epsilon\in(0,1)$. For every integer $K$ with
$(1-\epsilon)/\delta_K\le K\le(1-\epsilon/2)/\delta_K$,
\[
\norm{\Delta^N\big[\cU(t)[a\ue^{i\varphi_b/\hbar}]_\Gamma-u^{(K)}(t)_\Gamma\big]}\le C'\epsilon^{-1/2}C^{N+1}N^{2N}(\alpha\hbar)^{-2N}\ue^{-(1-\epsilon)\alpha^2/(C\hbar t)}\norm a ,
\]
with $C'$ an absolute constant.
\end{corollaryR}
The window for $K$ is the one used in Schubert's proof of his Corollary~1.5; the window printed in its statement
does not match that proof.

\begin{corollaryR}\label{cor:dd16}
If moreover the injectivity radius of $\Hb/\Gamma$ is at least $r_0>0$, then pointwise
\[
\big|\Delta^N\big[\cU(t)[a\ue^{i\varphi_b/\hbar}]_\Gamma-u^{(K)}(t)_\Gamma\big]\big|\le C^{N+K+1}(2N)!\,K!\,(\alpha\hbar)^{-2N-4}\Big(\frac{\hbar t}{\alpha^2}\Big)^{K+1}\norm a ,
\]
with $C$ depending also on $r_0$, and correspondingly for the window of Corollary~\ref{cor:dd15}.
\end{corollaryR}

\medskip\noindent\textbf{Notation for the analytic estimates.}
Let $X=y\partial_x$, $Y=y\partial_y$, $\Delta=X^2+Y^2-Y$, $\Delta(s)=\ue^{-2s}X^2+Y^2-Y$,
$\Lambda=\sqrt{-\Delta}$, $X_0=\cosh d(x_0,\cdot)$, $\tilde\psi=\log X_0$, and $W_\beta=X_0^\beta$.
In the analytic estimates below, the unadorned amplitude norm uses the weight $X_0^\beta$:
\[
\norm{a}:=\norm{\ue^{\alpha\Lambda}X_0^\beta a}.
\]
By Lemma~\ref{lem:m2w} a finite norm with the weight $\ue^{\beta\langle d\rangle}$ at radius $\alpha$ gives a finite norm of this
kind at the radius $c_0\alpha$; Theorem~\ref{thm:dd14} is deduced from Theorem~\ref{thm:dd-appx} in this way at the end of
Section~\ref{subsec:dd-thm}. The two-time mollified propagator $V_\varepsilon(s_1,s_0)=V_\varepsilon(s_1)V_\varepsilon(s_0)^*$ is defined in
Appendix~\ref{app:kl}.

\subsection{Exact identities}\label{subsec:dd-id}

Let $V(t)$ be the unitary operator of \eqref{eq:factor}, let $V(t,s):=V(t)V(s)^{-1}$, the propagator of
$i\partial_tV=-\tfrac{\hbar}{2}\Delta(t)V$ (\cite[(1.22)]{S}), and let $c:=i\hbar/2$. By \eqref{eq:factor} (which is \cite[(1.23)]{S}),
\[
\cU(t)(a\,\ue^{i\varphi/\hbar})=\ue^{-it/(2\hbar)}\bigl(S(t)V(t)a\bigr)\ue^{i\varphi/\hbar},
\qquad (S(t)w)(x,y)=\ue^{-t/2}w(x,\ue^{t}y),
\]
with $S(t)$ unitary. Then
\[S(t)^{-1}XS(t)=\ue^{-t}X,\qquad S(t)^{-1}YS(t)=Y,\]
hence $S(t)^{-1}\Delta S(t)=\Delta(t)$.
Multiplication by $\ue^{i\varphi/\hbar}=y^{-i/\hbar}$ preserves every Sobolev space $H^k$, and so do $\cU(t)$ and $S(t)$.
By the factorization, so do $V(t)^{\pm1}$ and $V(t,s)$.

\begin{proposition}[Exact identities (D1), (D3), (D4)]\label{prop:dd-id}
\begin{enumerate}
\item[\textup{(D1)}] \emph{Dyson expansion.} Let $a$ be such that every word in $X,Y$ applied to $a$ lies in $L^2(\Hb)$, and put
\[
P_k(t)=\int_{t>s_1>\cdots>s_k>0}\Delta(s_1)\cdots\Delta(s_k)\,ds,\qquad a^{(K)}(t)=\sum_{k=0}^{K}c^kP_k(t)a .
\]
Then, with $\Sigma_{K+1}(t)=\{t>s_1>\cdots>s_{K+1}>0\}$,
\[
a-V(t)^*a^{(K)}(t)=c^{K+1}\int_{\Sigma_{K+1}(t)}V(s_1)^*\,\Delta(s_1)\cdots\Delta(s_{K+1})a\,ds .
\]
This $P_k$ is exactly the $P_k$ of \cite{S}
(the display before the semiclassical approximation of order $K$). Every amplitude $a$ of Theorem~\ref{thm:dd-appx} satisfies the
hypothesis, by Lemma~\ref{lem:sc}.
\item[\textup{(D3)}] \emph{Transport of $\Delta^N$} (the conjugated operator is that of \cite[Proposition~2.2 and the proof of Proposition~3.10, (3.31)--(3.34)]{S}). Put $\Delta_\varphi:=e^{-i\varphi/\hbar}\Delta e^{i\varphi/\hbar}$, so that
$\hbar^2\Delta_\varphi=\hbar^2\Delta-2i\hbar Y-1+i\hbar$. Put $\Delta_\varphi(s):=S(s)^{-1}\Delta_\varphi S(s)$, so that
$\hbar^2\Delta_\varphi(s)=\hbar^2\Delta(s)-2i\hbar Y-1+i\hbar$. Then for every $b$ such that every word in $X,Y$ applied to $b$ lies in $L^2(\Hb)$,
\[
\Delta_\varphi^N V(s)^*b=V(s)^*\Delta_\varphi(s)^Nb.
\]
\item[\textup{(D4)}] \emph{The error.} Let $E:=\Delta^N\bigl[\cU(t)(a\,\ue^{i\varphi/\hbar})_\Gamma-u^{(K)}(t)_\Gamma\bigr]$, with $u^{(K)}$ as in \cite{S}. If every word in $X,Y$ applied to $a$ lies in $L^2_\omega$, for a weight $\omega$ as in Section~\ref{subsec:dd-surf}, then
\[
\begin{gathered}
\norm{E}_{L^2(M)}\le\Bigl(\frac{\hbar}{2}\Bigr)^{K+1}\int_{\Sigma_{K+1}(t)}\bigl\|\mathcal V_{s_1}b'(s)\bigr\|_{L^2(M)}\,ds,\\
b'(s):=\Delta_\varphi(s_1)^N\Delta(s_1)\cdots\Delta(s_{K+1})a.
\end{gathered}
\]
Everything of high order is now inside $b'(s)$, which is a fixed word-polynomial in $X,Y$ applied to $a$.
\end{enumerate}
\end{proposition}

\begin{proof}
\medskip\noindent\textbf{(D1).}

$V(s)^*=V(0,s)$ and $\partial_sV(t,s)=-c\,V(t,s)\Delta(s)$ on $H^\infty$. So for $b$ with all words in $L^2$, $s\mapsto V(s)^*b$ is
differentiable in $L^2$ with derivative $-c\,V(s)^*\Delta(s)b$. (In integrated form this is Lemma~\ref{lem:surf-duh} for the trivial
group, the weight $\omega=1$ and the constant path $g(s)=b$: then $\mathcal V_sg=\ue^{i\varphi/\hbar}V(s)^*g$ by \eqref{eq:factor}; the proof
of that lemma does not use (D1).)

$P_k(s)a$ is a finite sum of continuous scalar functions of $s$ times fixed words applied to $a$; $\partial_sP_k(s)=\Delta(s)P_{k-1}(s)$
and $P_k(0)=0$ for $k\ge1$.

With $J_k(t):=\int_0^tV(s)^*\Delta(s)P_k(s)a\,ds$ this gives $V(t)^*P_k(t)a=J_{k-1}(t)-cJ_k(t)$ for $k\ge1$, and $V(t)^*a-a=-cJ_0(t)$.

Multiplying by $c^k$ and summing telescopes to $a-V(t)^*a^{(K)}(t)=c^{K+1}J_K(t)$. Expanding $P_K$ gives the integral over $\Sigma_{K+1}(t)$.

\medskip\noindent\textbf{(D3).}

For $c'\in H^2$, the commutation of the propagator with the Laplacian gives the following identity.

$\cU(s)$ commutes with $\Delta$, so $\cU(s)(e^{i\varphi/\hbar}\Delta_\varphi c')=\Delta\,\cU(s)(e^{i\varphi/\hbar}c')$.

By \cite[(1.23)]{S} this reads $S(s)V(s)\Delta_\varphi c'=\Delta_\varphi S(s)V(s)c'$, i.e. $V(s)\Delta_\varphi=\Delta_\varphi(s)V(s)$ on $H^2$.

Multiplying by $V(s)^*$ on both sides is legitimate because $V(s)^*$ preserves $H^2$. Iterating $N$ times gives the assertion, since $b$ lies in every $H^k$.

\medskip\noindent\textbf{(D4).}

By definition $u^{(K)}(t)=\ue^{-it/(2\hbar)}\ue^{i\varphi/\hbar}S(t)a^{(K)}(t)$, so $u^{(K)}(t)_\Gamma=\cU(t)\mathcal V_ta^{(K)}(t)$. Since $\cU(t)$ is unitary on $L^2(M)$ and commutes with $\Delta_M$, $\norm E=\norm{\Delta_M^N(\mathcal V_0a-\mathcal V_ta^{(K)}(t))}$. (No commutation of $\cU(t)$ with periodization is used.)

By Corollary~\ref{cor:surf-dyson}, $\mathcal V_0a-\mathcal V_ta^{(K)}(t)=c^{K+1}\int_{\Sigma_{K+1}(t)}\mathcal V_{s_1}(\Delta(s_1)\cdots\Delta(s_{K+1})a)\,ds$.

By Lemma~\ref{lem:surf-lap} $N$ times, $\Delta_M^N\mathcal V_{s_1}g=\mathcal V_{s_1}\Delta_\varphi(s_1)^Ng$. Closedness of $\Delta_M^N$ and Minkowski's inequality give the bound, with $|c|^{K+1}=(\hbar/2)^{K+1}$.
(D1) and (D3) are identities on $\Hb$; on $M$ the proof of (D4) uses their surface versions in Section~\ref{subsec:dd-surf}. \qedhere
\end{proof}

\subsection{The surface object}\label{subsec:dd-surf}

Periodization needs decay, and Proposition~\ref{prop:ws} gives weighted bounds only for $V_\varepsilon$, not for $V$. So $[\ue^{i\varphi/\hbar}V(s)^*g]_\Gamma$ cannot be read as a convergent sum, and the identity $\cU(t)u_\Gamma=\sum_\gamma(\cU(t)u)\circ\gamma^{-1}$ of \cite{S} is not available for $u=\ue^{i\varphi/\hbar}V(s)^*g$. We therefore define this object on $M$ through the surface propagator, and prove the Duhamel formula on $M$. All estimates of the proofs below are unchanged.

Let $\omega\ge1$ be a locally bounded measurable weight on $\Hb$, $L^2_\omega:=\{f:\omega f\in L^2(\Hb)\}$, and $\delta_sz:=x+i\ue^sy$. Assume:
\begin{itemize}
\item[(W1)] $f\mapsto f_\Gamma$ is bounded from $L^2_\omega$ to $L^2(M)$, the sum converging absolutely a.e.;
\item[(W2)] $\sup\{\omega(z)/\omega(\delta_sz):|s|\le T,\ z\in\Hb\}<\infty$ for every $T$.
\end{itemize}
Lemma~\ref{lem:t2} (its proof, $t=0$) gives (W1) for $\omega=W_\sigma$, $\sigma>1/2$. Proposition~\ref{os:p1} gives it for the periodization weight $\omega=X_0^{1/2}\ell^{\,c'}$, $\ell=\log(2X_0)/\log2$, of Section~\ref{sec:printed}. Since $d(z,\delta_sz)=|s|$ and $\tilde\psi$ is $1$-Lipschitz, (W2) holds with the explicit ratios $W_\sigma(z)/W_\sigma(\delta_sz)\le\ue^{\sigma|s|}$ and $\omega(z)/\omega(\delta_sz)\le\ue^{|s|/2}(1+|s|/\log 2)^{c'}$. By (W2), $S(s)$ is bounded on $L^2_\omega$, locally uniformly in $s$; it is strongly continuous there, by density of $C_c(\Hb)$.

Let $\mathcal D_\omega$ be the set of $g\in L^2_\omega$ with $Xg,Yg,X^2g,Y^2g\in L^2_\omega$ (weak derivatives). For $g\in L^2_\omega$ and $s\in\mathbb R$ put
\[
\mathcal V_sg:=\ue^{-is/(2\hbar)}\,\cU(-s)\,\bigl[\ue^{i\varphi/\hbar}S(s)g\bigr]_\Gamma\in L^2(M),
\]
with $\cU$ the propagator on $M$. Thus $\mathcal V_0g=[\ue^{i\varphi/\hbar}g]_\Gamma$ and $\norm{\mathcal V_sg}_{L^2(M)}=\norm{[\ue^{i\varphi/\hbar}S(s)g]_\Gamma}_{L^2(M)}$. This $\mathcal V_sg$ is what $[\ue^{i\varphi/\hbar}V(s)^*g]_\Gamma$ formally means by \cite[(1.23)]{S}.

\begin{lemma}[Laplacian on the weighted class]\label{lem:surf-lap}
For $g\in\mathcal D_\omega$ and $s\in\mathbb R$, $[\ue^{i\varphi/\hbar}S(s)g]_\Gamma\in\operatorname{Dom}(\Delta_M)$ and $\Delta_M[\ue^{i\varphi/\hbar}S(s)g]_\Gamma=[\ue^{i\varphi/\hbar}S(s)\Delta_\varphi(s)g]_\Gamma$. Hence $\mathcal V_sg\in\operatorname{Dom}(\Delta_M)$ and $\Delta_M\mathcal V_sg=\mathcal V_s\Delta_\varphi(s)g$.
\end{lemma}

\begin{proof}
\begin{itemize}
\item Put $f:=\ue^{i\varphi/\hbar}S(s)g$. Then $Xf=\ue^{-s}\ue^{i\varphi/\hbar}S(s)Xg$, $Yf=\ue^{i\varphi/\hbar}S(s)(Y-i/\hbar)g$ and $\Delta f=\ue^{i\varphi/\hbar}S(s)\Delta_\varphi(s)g$. So $f,Xf,Yf,\Delta f\in L^2_\omega$ by (W2), and $f\in\operatorname{Dom}(\Delta)$ because $\Delta$ is essentially self-adjoint on $C_c^\infty(\Hb)$.
\item Take radial $\chi_n=\rho(d(x_0,\cdot)/n)$, with $\rho\in C_c^\infty$, $\rho=1$ near $0$. Then $0\le\chi_n\le1$, $\chi_n\to1$, and $|X\chi_n|+|Y\chi_n|+|\Delta\chi_n|\le C/n$. Since $\Delta(\chi_nf)=\chi_n\Delta f+2(X\chi_n\,Xf+Y\chi_n\,Yf)+(\Delta\chi_n)f$, dominated convergence gives $\chi_nf\to f$ and $\Delta(\chi_nf)\to\Delta f$ in $L^2_\omega$.
\item Fix $n$. Essential self-adjointness gives test functions $F_m\to\chi_nf$ in graph norm. With a fixed $\theta\in C_c^\infty(\Hb)$, $\theta=1$ on $\operatorname{supp}\chi_n$, the test functions $\theta F_m$ converge to $\chi_nf$ in graph norm, since $\norm{\nabla F}^2=\langle F,-\Delta F\rangle$ controls the cross term. Their supports lie in the compact set $\operatorname{supp}\theta$, where $\omega$ is bounded.
\item For a test function $F$, $F_\Gamma\in\operatorname{Dom}(\Delta_M)$ and $\Delta_MF_\Gamma=(\Delta F)_\Gamma$. By (W1) and closedness of $\Delta_M$, first in $m$ and then in $n$, the claim for $f$ follows. The statement for $\mathcal V_s$ follows since $\cU(-s)$ commutes with $\Delta_M$.\qedhere
\end{itemize}
\end{proof}

\begin{lemma}[Duhamel on $M$]\label{lem:surf-duh}
Let $g\in C^1([0,T];L^2_\omega)$ with $g(s)\in\mathcal D_\omega$ and $Yg,X^2g,Y^2g\in C([0,T];L^2_\omega)$. Then, with $c=i\hbar/2$,
\[
\mathcal V_Tg(T)-\mathcal V_0g(0)=\int_0^T\mathcal V_s\bigl(g'(s)-c\,\Delta(s)g(s)\bigr)\,ds\qquad\text{in }L^2(M).
\]
\end{lemma}

\begin{proof}
\begin{itemize}
\item For $g_0$ with $g_0,Yg_0\in L^2_\omega$: $S(s)g_0-S(r)g_0=\int_r^sS(\tau)(Y-\tfrac12)g_0\,d\tau$. This holds in $L^2(\Hb)$ (generator of the dilation group), the right side converges in $L^2_\omega$ by strong continuity, and $L^2_\omega\subset L^2$ since $\omega\ge1$.
\item Hence $s\mapsto S(s)g(s)$ is $C^1$ into $L^2_\omega$, and by (W1) $h(s):=[\ue^{i\varphi/\hbar}S(s)g(s)]_\Gamma$ is $C^1$ into $L^2(M)$ with $h'=[\ue^{i\varphi/\hbar}S(s)((Y-\tfrac12)g+g')]_\Gamma$.
\item By Lemma~\ref{lem:surf-lap}, $h(s)\in\operatorname{Dom}(\Delta_M)$ and $\Delta_Mh=[\ue^{i\varphi/\hbar}S(s)\Delta_\varphi(s)g]_\Gamma$, continuous in $s$. So $\tfrac{d}{ds}\cU(-s)h(s)=\cU(-s)\bigl(h'(s)-\tfrac{i\hbar}{2}\Delta_Mh(s)\bigr)$.
\item Since $\Delta_\varphi(s)=\Delta(s)-\tfrac{2i}{\hbar}(Y-\tfrac12)-\hbar^{-2}$, we have $-\tfrac{i\hbar}{2}\Delta_\varphi(s)=-c\Delta(s)-(Y-\tfrac12)+\tfrac{i}{2\hbar}$. The $(Y-\tfrac12)$ terms cancel, and $\tfrac{i}{2\hbar}$ cancels the derivative of $\ue^{-is/(2\hbar)}$. Thus $\tfrac{d}{ds}\mathcal V_sg(s)=\mathcal V_s(g'(s)-c\Delta(s)g(s))$; integrate.\qedhere
\end{itemize}
\end{proof}

\begin{corollary}[Dyson expansion on $M$]\label{cor:surf-dyson}
If every word in $X,Y$ applied to $a$ lies in $L^2_\omega$, then with $a^{(K)}(t)=\sum_{k\le K}c^kP_k(t)a$,
\[
\mathcal V_0a-\mathcal V_ta^{(K)}(t)=c^{K+1}\int_{\Sigma_{K+1}(t)}\mathcal V_{s_1}\bigl(\Delta(s_1)\cdots\Delta(s_{K+1})a\bigr)\,ds.
\]
\end{corollary}

\begin{proof}
$P_k(s)a$ is a finite sum of continuous scalar functions times fixed words applied to $a$, so Lemma~\ref{lem:surf-duh} applies to it. Since $\partial_sP_{k}(s)=\Delta(s)P_{k-1}(s)$ for $k\ge1$ and $P_k(0)=0$, with $J_k(t):=\int_0^t\mathcal V_s\Delta(s)P_k(s)a\,ds$ we get $\mathcal V_tP_k(t)a=J_{k-1}(t)-cJ_k(t)$ for $k\ge1$, and $\mathcal V_ta-\mathcal V_0a=-cJ_0(t)$. Multiplying by $c^k$ and summing telescopes to $\mathcal V_0a-\mathcal V_ta^{(K)}(t)=c^{K+1}J_K(t)$. Expanding $P_K$ gives the integral over $\Sigma_{K+1}(t)$.
\end{proof}

\subsection{First-order Cauchy estimates (Lemma FC)}\label{subsec:dd-fc}

\begin{lemma}[FC]\label{lem:fc}
For $\tau_0<\pi/2$ there is $C_1=C_1(\tau_0)$ such that for $Z\in\{X,Y\}$ and $0\le\alpha''<\alpha'\le\tau_0$,
\[
\norm{\ue^{\alpha''\Lambda}Z\,\ue^{-\alpha'\Lambda}}_{L^2(\Hb)}\le\frac{C_1}{\alpha'-\alpha''}.
\]
\end{lemma}

\begin{proof}
\medskip\noindent\textbf{Reduction.}

We first Fourier-transform in $x$. On the mode $\ue^{ikx}$, $k\ne0$, put $z=|k|y$ and write the function of $y$
as $z^{1/2}f(z)$, with $f\in L^2(dz/z)$. Then $-\Delta-\tfrac14=-(z\partial_z)^2+z^2$,
$X=i\,\operatorname{sgn}(k)\,z$ is multiplication, and $Y=z\partial_z+\tfrac12$.
Thus $Y_0:=Y-\tfrac12=z\partial_z=:D_z$ is skew-adjoint on $L^2(dz/z)$.

In the unitary Kontorovich--Lebedev picture
\[G=w(\nu)^{1/2}F,\qquad F(\nu)=\int_0^\infty K_{i\nu}(z)f(z)\,dz/z,\qquad w(\nu)=\tfrac{2}{\pi^2}\nu\sinh\pi\nu\]
(\cite{Y}), $\Lambda$ acts as multiplication by $\Lambda(\nu)=(\nu^2+\tfrac14)^{1/2}$.

The operators are the same for every $k$, so bounds are uniform in $k$. On the $k=0$ mode, $X=0$, while $Y_0$ and $\Lambda$ commute: in $r=\log y$, $Y_0=\partial_r$ and $\Lambda=(-\partial_r^2+\tfrac14)^{1/2}$. So
\[\norm{\ue^{\alpha''\Lambda}Y_0\ue^{-\alpha'\Lambda}}\le\norm{\Lambda\ue^{-(\alpha'-\alpha'')\Lambda}}\]
there. The mode $k=0$ has measure zero in the direct integral.

Put $\gamma:=\alpha'-\alpha''$ and write
\[
\ue^{\alpha''\Lambda}Z\ue^{-\alpha'\Lambda}=\bigl(\ue^{\alpha''\Lambda}Z\ue^{-\alpha''\Lambda}\Lambda^{-1}\bigr)\bigl(\Lambda\ue^{-\gamma\Lambda}\bigr).
\]
The second factor has norm $\le\sup_{\lambda\ge1/2}\lambda\ue^{-\gamma\lambda}\le1/(e\gamma)$. It remains to bound the first factor uniformly for $\alpha''\le\tau_0$, and $Y=Y_0+\tfrac12$ adds at most $\tfrac12$.

\medskip\noindent\textbf{Spectral notation.} Put $p:=\pi(\nu+\nu')/2$, $q:=\pi(\nu-\nu')/2$, $x:=\nu-\nu'$. Since $\sinh\pi\nu\,\sinh\pi\nu'=\sinh(p+q)\sinh(p-q)$,
\[
(w(\nu)w(\nu'))^{1/2}=\tfrac{2}{\pi^2}(\nu\nu')^{1/2}(\sinh^2p-\sinh^2q)^{1/2}=\tfrac{2}{\pi^2}(\nu\nu')^{1/2}(\cosh^2p-\cosh^2q)^{1/2}.
\]

\medskip\noindent\textbf{Mellin product formula.} For $\operatorname{Re}s>|\operatorname{Re}a|+|\operatorname{Re}b|$,
\[
M_s(a,b):=\int_0^\infty z^{s-1}K_a(z)K_b(z)\,dz=2^{s-3}\,\frac{\Gamma\bigl(\tfrac{s+a+b}{2}\bigr)\Gamma\bigl(\tfrac{s+a-b}{2}\bigr)\Gamma\bigl(\tfrac{s-a+b}{2}\bigr)\Gamma\bigl(\tfrac{s-a-b}{2}\bigr)}{\Gamma(s)}.
\]
\emph{Proof of the formula.} Write $K_b=K_{-b}$ (\cite[10.27.3]{DLMF}) and use the integral representation
\cite[10.32.10]{DLMF} for both Bessel factors. The integration in $z$ is
\[
\int z^{m-1}\ue^{-z^2c/4}\,dz=\tfrac12(4/c)^{m/2}\Gamma(m/2),
\qquad m=s+a-b,\quad c=1/u+1/v.
\]
Substitute $u=rt$ and $v=r(1-t)$ in the remaining double integral. The $r$-integral gives
$\Gamma(\tfrac{s-a+b}{2})$, and the $t$-integral gives the beta function
$\Gamma(\tfrac{s-a-b}{2})\Gamma(\tfrac{s+a+b}{2})/\Gamma(s)$, yielding the formula above.
The absolute integrand is the same expression with $a,b$ replaced by $\operatorname{Re}a,\operatorname{Re}b$ (since $|(z/2)^a|=(z/2)^{\operatorname{Re}a}$). For $\operatorname{Re}s>|\operatorname{Re}a|+|\operatorname{Re}b|$ the same computation then gives a finite value with all Gamma arguments of positive real part, so Tonelli justifies every exchange. This covers $a=i\nu$, $b=i\nu'$ at $s=1$, and $b=i\nu'\pm1$ at $s=1+\eta$, $\eta>0$. \qed

\medskip\noindent\textbf{(X) Kernel of $z$.}

Since the measure is $dz/z$, the kernel of multiplication by $z$ is $(w w')^{1/2}M_1(i\nu,i\nu')$.

The identity $|\Gamma(\tfrac12+iy)|^2=\pi/\cosh\pi y$ gives
\[
Zt_1(\nu,\nu')=(ww')^{1/2}\frac{\pi^2}{4\cosh p\cosh q}=\frac{(\nu\nu')^{1/2}(\cosh^2p-\cosh^2q)^{1/2}}{2\cosh p\cosh q},
\]
so $|Zt_1|\le(\nu\nu')^{1/2}/(2\cosh q)$.

The kernel of $\ue^{\alpha''\Lambda}z\ue^{-\alpha''\Lambda}\Lambda^{-1}$ is bounded by $\ue^{\alpha''|x|}(\nu\nu')^{1/2}/(2\cosh(\pi x/2)\Lambda(\nu'))$.

Since $(\nu\nu')^{1/2}\le(\nu+\nu')/2\le\nu'+|x|/2$ and $\Lambda(\nu')\ge\max(\nu',1/2)$, we have
\[
(\nu\nu')^{1/2}/\Lambda(\nu')\le1+|x|.
\]

So the kernel is at most $g(x):=(1+|x|)\ue^{\tau_0|x|}/(2\cosh(\pi x/2))$, which is integrable for $\tau_0<\pi/2$. The Schur test gives the bound $\int g$.

\medskip\noindent\textbf{(Y) Kernel of $D_z$.} The integral $\int K_{i\nu}K'_{i\nu'}\,dz$ does not converge at $z=0$, so we regularize.

Let $F,G\in C_c^\infty((0,\infty))$ be the KL transforms of $f,g$ (the core $\mathcal C$).

On $\mathcal C$, $f(e^{-L})$ and $zg'(z)|_{z=e^{-L}}$ decay faster than any power of $L$ as $L\to\infty$: they are integrals of smooth compactly supported densities against $\ue^{\pm i\nu L}$-oscillations, plus $O(\ue^{-2L})$. This uses $K_{i\nu}=\tfrac12\pi(I_{-i\nu}-I_{i\nu})/\sin(i\nu\pi)$ (\cite[10.27.4]{DLMF}) with the series of $I_{\pm i\nu}$ (\cite[10.25.2]{DLMF}), and $1/\sin(i\nu\pi)$ is smooth on $\operatorname{supp}F\subset(0,\infty)$.

Hence $\langle f,D_zg\rangle=\int_0^\infty\bar f g'\,dz=\lim_{\eta\to0^+}\int z^\eta\bar fg'\,dz$, by dominated convergence.

For $\eta>0$ Fubini applies, and $K'_\nu=-(K_{\nu-1}+K_{\nu+1})/2$ (\cite[10.29.1]{DLMF}, with $\mathcal Z_\nu=\ue^{\nu\pi i}K_\nu$ in DLMF's convention, under which 10.29.1 reads $\mathcal Z_{\nu-1}+\mathcal Z_{\nu+1}=2\mathcal Z'_\nu$ and gives the minus sign for $K$) gives
\[
\begin{gathered}
\langle f,D_zg\rangle=\lim_{\eta\to0^+}\iint\bar F(\nu)G(\nu')\,w(\nu)w(\nu')\,I_\eta(\nu,\nu')\,d\nu\,d\nu',\\
I_\eta=-\tfrac12\bigl[M_{1+\eta}(i\nu,i\nu'+1)+M_{1+\eta}(i\nu,i\nu'-1)\bigr].
\end{gathered}
\]
Both terms are admissible, since $\operatorname{Re}s=1+\eta>0+1$.

In $M_{1+\eta}(i\nu,i\nu'+1)$ the factor $\Gamma((\eta+ix)/2)=(2/(\eta+ix))\Gamma(1+(\eta+ix)/2)$ is singular, and the other factors are smooth. In $M_{1+\eta}(i\nu,i\nu'-1)$ the singular factor is $\Gamma((\eta-ix)/2)$.

So $I_\eta=\frac{2}{\eta+ix}A_\eta+\frac{2}{\eta-ix}B_\eta$, where $A_\eta,B_\eta$ are smooth in $(\nu,\nu',\eta)$ on the support.

At $\eta=0$ and $x=0$, the reflection identity $\Gamma(1+w)\Gamma(-w)=-\pi/\sin\pi w$ gives
\[
\begin{aligned}
A_0&=-\tfrac12\cdot\tfrac14\,\Gamma(1+ip/\pi)\Gamma(-ip/\pi)=-\tfrac18\cdot\tfrac{i\pi}{\sinh p},\\
B_0&=-\tfrac12\cdot\tfrac14\,\Gamma(ip/\pi)\Gamma(1-ip/\pi)=-\tfrac18\cdot\tfrac{-i\pi}{\sinh p}.
\end{aligned}
\]
So $A_0+B_0=0$ on the diagonal.

The Sokhotski--Plemelj formula gives $\frac{2}{\eta\pm ix}\to\mp2i\,\mathrm{PV}\frac1x+2\pi\delta(x)$. The $\delta$-terms carry $A_0+B_0$ at $x=0$, so they cancel.

For $x\neq0$ the pointwise limit is $I_0=-\pi^2/(4\sinh p\,\sinh q)$ (same Gamma reflections). Hence, on $\mathcal C\times\mathcal C$,
\[
\langle f,D_zg\rangle=\mathrm{PV}\iint\bar F G\,w\,w'\,I_0,
\qquad\text{i.e.}\qquad
\begin{gathered}
\tilde D(\nu,\nu')=-\frac{(\nu\nu')^{1/2}\rho(\nu,\nu')}{2\sinh q},\\
\rho:=\Bigl(1-\frac{\sinh^2q}{\sinh^2p}\Bigr)^{1/2}\in[0,1].
\end{gathered}
\]

This kernel is real and antisymmetric, as $D_z$ must be.

As a consistency check, not used in the proof, $(\nu^2-\nu'^2)\tilde D=-2Zt_2$, where
$Zt_2=(w(\nu)w(\nu'))^{1/2}M_2(\nu,\nu')$ is the kernel of $z^2$, with $M_2$ as in
Remark~\ref{rem:fc-y-alt}. This agrees with $[\Lambda^2,D_z]=[z^2,D_z]=-2z^2$.

\medskip\noindent\textbf{(Y) bound.} The kernel of $Q:=\ue^{\alpha''\Lambda}D_z\ue^{-\alpha''\Lambda}\Lambda^{-1}$ is $\ue^{\alpha''(\Lambda(\nu)-\Lambda(\nu'))}\tilde D(\nu,\nu')/\Lambda(\nu')$. It is valid on $\mathcal C$, because $\ue^{\pm\alpha''\Lambda}$ and $\Lambda^{-1}$ preserve $\mathcal C$. Split $Q=Q_{\rm far}+Q_{\rm near}$, with $Q_{\rm near}:=Q\cdot\mathbf 1_{|x|\le1}$.

\textbf{$Q_{\rm far}$.} Here $|Q|\le(1+|x|)\ue^{\tau_0|x|}/(2\sinh(\pi|x|/2))$ on $|x|>1$. This is integrable, so the Schur test applies.

\textbf{$Q_{\rm near}$.} Write $1/(2\sinh(\pi x/2))=1/(\pi x)+r(x)$, with $r$ bounded on $|x|\le1$.

The $r$-part has a bounded kernel supported in $|x|\le1$, so the Schur test applies.

In the main part $-a(\nu,\nu')/(\pi(\nu-\nu'))$, put
\[a(\nu,\nu'):=\ue^{\alpha''(\Lambda(\nu)-\Lambda(\nu'))}(\nu\nu')^{1/2}\rho(\nu,\nu')/\Lambda(\nu')\]
Then $a(\nu,\nu')=a(\nu',\nu')+[a(\nu,\nu')-a(\nu',\nu')]$, with $a(\nu',\nu')=\nu'/\Lambda(\nu')\in[0,1]$.

The diagonal part is the truncated Hilbert transform $\mathrm{PV}\,\frac{1}{\pi x}\mathbf 1_{|x|\le1}$
composed with multiplication by $\nu'/\Lambda(\nu')$. This transform is bounded on $L^2(\mathbb R)$:
its Fourier multiplier $-\tfrac{2i}{\pi}\operatorname{Si}(|\xi|)\operatorname{sgn}\xi$ has modulus
at most $(2/\pi)\operatorname{Si}(\pi)<1.2$. Extending functions by zero from $(0,\infty)$ therefore gives the required bound.

For the off-diagonal remainder, write $s(\nu):=(\nu\nu')^{1/2}/\Lambda(\nu')$ and
$E:=\ue^{\alpha''(\Lambda(\nu)-\Lambda(\nu'))}$. Then
\[
a(\nu,\nu')-a(\nu',\nu')=(E-1)s\rho+(\rho-1)s+(s(\nu)-s(\nu')).
\]
For $|x|\le1$, the first term is controlled by $|E-1|\le\tau_0\ue^{\tau_0}|x|$ and
$s\rho\le s\le1+|x|\le2$. For the second, $\sinh|q|/\sinh p\le|q|/p$ when $|q|\le p$ gives
$1-\rho\le\sinh^2q/\sinh^2p\le|q|/p=|x|/(\nu+\nu')$, and hence
$(1-\rho)s\le|x|(\nu\nu')^{1/2}/((\nu+\nu')\Lambda(\nu'))\le|x|$.
Finally, $|s(\nu)-s(\nu')|=|\sqrt\nu-\sqrt{\nu'}|\sqrt{\nu'}/\Lambda(\nu')\le|x|^{1/2}$,
since $\sqrt{\nu'}/\Lambda(\nu')\le1$.
So the remainder kernel is $\le C|x|^{-1/2}$ on $|x|\le1$, which is integrable. The Schur test therefore bounds the remainder.
Combining the near and far bounds gives $\norm{Q}\le C(\tau_0)$ on $\mathcal C$, so $Q$ extends to a bounded operator. The bound for $\ue^{\alpha''\Lambda}D_z\ue^{-\alpha'\Lambda}$ follows on the dense set $\Lambda\ue^{-\gamma\Lambda}\mathcal C$, and then by closure. \qed
\end{proof}

\begin{remark}[Proof of (Y) without a principal value]\label{rem:fc-y-alt}
The case $Z=Y$ of Lemma~\ref{lem:fc} also follows without the kernel of $D_z$. For $A,B\in C_c^\infty((0,\infty))$ put $a=wA$, $b=wB$, $\hat a(t)=\int a(\nu)\ue^{i\nu t}\,d\nu$, $\psi_a(z)=\int a(\nu)K_{i\nu}(z)\,d\nu=\tfrac12\int\ue^{-z\cosh t}\hat a(t)\,dt$ and $S(A,B)=\langle\psi_a,D_z\psi_b\rangle_{L^2(dz/z)}$.
\begin{itemize}
\item By Fubini, $S(A,B)=-\tfrac14\iint\overline{\hat a(t)}\,\hat b(u)\,\frac{\cosh u}{\cosh t+\cosh u}\,dt\,du$, absolutely convergent.
\item Since $(\partial_u^2-\partial_t^2)\frac{\cosh u}{\cosh t+\cosh u}=2(\cosh t+\cosh u)^{-2}$, integrating by parts twice in each variable gives
\[S(\Lambda^2A,B)-S(A,\Lambda^2B)=-2\iint\bar A(\nu)B(\nu')\,w(\nu)w(\nu')\,M_2(\nu,\nu')\,d\nu\,d\nu',\]
where $M_2(\nu,\nu')=\int_0^\infty zK_{i\nu}K_{i\nu'}\,dz=\tfrac12\,\mathsf g(\tfrac{\nu+\nu'}2)\,\mathsf g(\tfrac{\nu-\nu'}2)$ and $\mathsf g(y)=\pi y/\sinh\pi y$ (the Mellin product formula at $s=2$).
\item The multipliers $\ue^{\pm s\Lambda}$ preserve $C_c^\infty((0,\infty))$. Integrating first along $\ue^{-t\Lambda}$ in both arguments over $t\in(0,\infty)$, then along $(\ue^{s\Lambda},\ue^{-s\Lambda})$ over $s\in(0,a)$, gives
\[S(\ue^{a\Lambda}A,\ue^{-a\Lambda}B)-S(A,B)=\iint\bar A(\nu)B(\nu')\,w(\nu)w(\nu')\,R_a(\nu,\nu')\,d\nu\,d\nu',\]
with $R_a=-2M_2\,\phi_a(\Lambda(\nu)-\Lambda(\nu'))/(\Lambda(\nu)+\Lambda(\nu'))$ and $\phi_a(\delta)=\int_0^a\ue^{s\delta}\,ds$.
\item Hence, on products in the KL picture and then by density, $\ue^{a\Lambda}Y\ue^{-(a+\gamma)\Lambda}=Y\ue^{-\gamma\Lambda}+T_{R_a}\ue^{-\gamma\Lambda}$, where $T_RG(\nu)=\int R(\nu,\nu')G(\nu')w(\nu')\,d\nu'$. The first term has norm at most $2/\gamma$, by $\norm{Yu}^2\le\norm{u}\,\norm{\Delta u}$. For the second, $\mathsf g(p')\cosh(\pi p')\le1+\pi p'$ with $p'=(\nu+\nu')/2$, together with the inequalities used in (X), gives
\[(w(\nu)w(\nu'))^{1/2}\,|R_a(\nu,\nu')|\,\ue^{-\gamma\Lambda(\nu')}\le C(\tau_0)\,\gamma^{-1}(1+|x|)\ue^{\tau_0|x|}\,\mathsf g(x/2),\qquad x=\nu-\nu',\]
which is integrable for $\tau_0<\pi/2$, and the Schur test concludes.
\end{itemize}
\end{remark}

\subsection{Unweighted Cauchy estimate for words (Lemma WC)}\label{subsec:dd-wc}

\begin{lemma}[WC]\label{lem:wc}
For a word $Z_1\cdots Z_m$ with $Z_i\in\{X,Y\}$, and $0\le\alpha''<\alpha\le\tau_0<\pi/2$:
\[
\norm{\ue^{\alpha''\Lambda}Z_1\cdots Z_m\ue^{-\alpha\Lambda}}\le\Bigl(\frac{C_1m}{\alpha-\alpha''}\Bigr)^m\le\Bigl(\frac{eC_1}{\alpha-\alpha''}\Bigr)^m m!.
\]
\end{lemma}

\begin{proof}
Take equally spaced radii $\alpha''=r_0<r_1<\cdots<r_m=\alpha$, and insert $\ue^{-r_i\Lambda}\ue^{r_i\Lambda}$.
Each factor $\ue^{r_{i-1}\Lambda}Z_i\ue^{-r_i\Lambda}$ is bounded by $C_1m/(\alpha-\alpha'')$ (Lemma~\ref{lem:fc}). The identity extends from $\mathcal C$ by closure, since each $\ue^{r_i\Lambda}$ is closed. Finally $m^m\le e^m m!$. \qed
\end{proof}

\subsection{Static weighted Cauchy estimates (Lemma SC)}\label{subsec:dd-sc}

Let $\mathcal A:=\operatorname{span}\{y,\;x,\;(x^2+1)/y\}$.
\begin{itemize}
\item $\mathcal A$ contains $X_0=\tfrac12\bigl((x^2+1)/y+y\bigr)$.
\item $\mathcal A$ is invariant under $X$ and $Y$: $Xy=0$, $Xx=y$, $X((x^2+1)/y)=2x$; $Yy=y$, $Yx=0$, $Y((x^2+1)/y)=-(x^2+1)/y$.
\item In this basis the $\ell^1$-operator norms are $\norm{X}\le2$ and $\norm{Y}\le1$.
\item Since $X_0\ge y/2$, $X_0\ge|x|$ (because $x^2+y^2+1-2y|x|=(|x|-y)^2+1$), and $X_0\ge(x^2+1)/(2y)$, every $v=c_1y+c_2x+c_3(x^2+1)/y$ satisfies $|v|\le2\norm{c}_1X_0$.
\item Hence $|Z_BX_0|\le2\cdot2^{|B|}X_0$ for every word $Z_B$.
\end{itemize}

\begin{lemma}[Faà di Bruno bound]\label{lem:fdb}
For $\beta\ge0$ and a word $Z_S$ of length $j$:
\[
|X_0^\beta Z_SX_0^{-\beta}|\le B_\beta\,j!\,\rho_\beta^{-j},\qquad \rho_\beta:=\tfrac14\log(3/2)
\]
(any fixed value below $\tfrac12\log(3/2)$ works), where $B_\beta$ depends only on $\beta$.
\end{lemma}

\begin{proof}
\begin{itemize}
\item Expand by induction on $j$: $Z_S\varphi(X_0)=\sum_{\pi}\varphi^{(|\pi|)}(X_0)\prod_{B\in\pi}Z_BX_0$, the sum over partitions $\pi$ of $S$, where each block $B$ inherits the order of $S$.
\item With $\varphi(u)=u^{-\beta}$, $|\varphi^{(k)}(u)|=(\beta)_ku^{-\beta-k}$. So the bound is $2^j\sum_kS(j,k)(\beta)_k2^k$.
\item The exponential generating function of $\sum_kS(j,k)(\beta)_k2^k$ is $(3-2\ue^z)^{-\beta}$, which is holomorphic on $|z|<\log(3/2)$.
\item Cauchy's estimate on $|z|=\rho<\log(3/2)$ gives the bound $C_{\beta,\rho}\,j!\,\rho^{-j}$. Combined with the factor $2^j$, this gives $\rho_\beta=\rho/2$. \qed
\end{itemize}
\end{proof}

\begin{lemma}[SC]\label{lem:sc}
For $\beta\ge0$, a word $Z_1\cdots Z_m$ and $0<\alpha\le\tau_0$, with $F:=X_0^\beta a$ and $\norm{\ue^{\alpha\Lambda}F}<\infty$:
\[
\norm{X_0^\beta Z_1\cdots Z_m a}\le B_\beta(m+1)\,m!\,(C_3/\alpha)^m\,\norm{\ue^{\alpha\Lambda}X_0^\beta a},
\qquad C_3:=\max(\ue C_1,\tau_0/\rho_\beta).
\]
\end{lemma}

\begin{proof}
\begin{itemize}
\item Leibniz for words of derivations: \[X_0^\beta Z_1\cdots Z_m(X_0^{-\beta}F)=\sum_{S\subseteq[m]}(X_0^\beta Z_SX_0^{-\beta})\,(Z_{S^c}F)\]
pointwise.
\item Bound each term by $B_\beta|S|!\rho_\beta^{-|S|}\norm{Z_{S^c}F}$, and the second factor by Lemma~\ref{lem:wc} with $\alpha''=0$: $\norm{Z_{S^c}F}\le(\ue C_1/\alpha)^{m-|S|}(m-|S|)!\,\norm{\ue^{\alpha\Lambda}F}$.
\item Since $\sum_{|S|=j}j!(m-j)!=m!$, the sum is $\le B_\beta\,m!\sum_j\rho_\beta^{-j}(\ue C_1/\alpha)^{m-j}\le B_\beta(m+1)m!(C_3/\alpha)^m$. \qed
\end{itemize}
\end{proof}

\subsection{\texorpdfstring{Bounds for $b'(s)$ (Lemma WB)}{Bounds for b'(s) (Lemma WB)}}\label{subsec:dd-wb}

\begin{lemma}[WB]\label{lem:wb}
Assume $0<\hbar\le1$ and $0<\alpha\le\alpha_0\le\min(1,\tau_0)$. Let $J$ be a word of length $\le n$, and let $\mathfrak D^J:=Z_1\cdots Z_m$ with $Z_i\in\{X,Y\}$ and $m=|J|$ be the ordered word it determines (the empty word is the identity). Uniformly in $s\in\Sigma_{K+1}(t)$:
\begin{itemize}
\item \textup{(WB-w)} $\norm{X_0^\beta\mathfrak D^Jb'(s)}\le B_\beta C_4^{N+K+n+1}(2N)!\,(n+2K+2)!\,\alpha^{-(n+2K+2)}(\alpha\hbar)^{-2N}\norm a$;
\item \textup{(WB-u)} for $0\le\alpha_1<\alpha_0'\le\alpha_0$ with $\alpha_0'-\alpha_1\ge\alpha/4$:
\[
\norm{\ue^{\alpha_1\Lambda}b'(s)}\le C_4^{N+K+1}(2N)!\,(2K+2)!\,\alpha^{-(2K+2)}(\alpha\hbar)^{-2N}\norm{\ue^{\alpha_0'\Lambda}a}.
\]
\end{itemize}
\end{lemma}

\begin{proof}
\begin{itemize}
\item Expand $b'$ as a sum over choices of one term from each factor:
  \begin{itemize}
  \item from each $\Delta(s_j)$: $\ue^{-2s_j}X^2$, $Y^2$ or $-Y$ (coefficients of modulus $\le1$, length $\le2$);
  \item from each $\Delta_\varphi(s_1)$: $\hbar^{-2}(-1+i\hbar)$ (length $0$, $|\text{coeff}|\le2\hbar^{-2}$), $-2i\hbar^{-1}Y$ (length $1$, $|\text{coeff}|=2\hbar^{-1}$), or one of the three length-$\le2$ terms (coefficient $\le1$).
  \end{itemize}
\item Suppose $n_0,n_1,n_2$ of the $N$ $\Delta_\varphi$-factors choose these three kinds. The coefficient is $\le2^{n_0+n_1}\hbar^{-2n_0-n_1}$, and the word length is $\le L:=n+\ell+2K+2$, with $\ell:=n_1+2n_2\le2N$.
\item Since $2n_0+n_1=2N-\ell$, the coefficient is $\le2^N\hbar^{-2N+\ell}$.
\item By Lemma~\ref{lem:sc}, and since $m\mapsto(m+1)m!(C_3/\alpha)^m$ is increasing, each term is $\le B_\beta(L+1)L!(C_3/\alpha)^L\norm a$.
\item Use $L!\le2^L\ell!(n+2K+2)!$, and then
\[
\begin{aligned}
\hbar^{-2N+\ell}\ell!(C_3/\alpha)^\ell&=\hbar^{-2N}\ell!(C_3\hbar/\alpha)^\ell\\
&\le\hbar^{-2N}(2N)!\max(1,C_3\hbar/\alpha)^{2N}\le(2N)!\,C_3^{2N}(\alpha\hbar)^{-2N},
\end{aligned}
\]
using $\alpha,\hbar\le1\le C_3$.
\item The per-factor weights for the $\Delta_\varphi$-factors are $2+2+1+1+1=7$ (weight $2$ for the length-$0$ and length-$1$ choices, weight $1$ for each of the three length-$\le2$ choices), so the weighted number of choices is $\le7^N3^{K+1}$, the factor $3^{K+1}$ coming from the $\Delta(s_j)$. Collecting the powers of $2$ and $7$ into $C_4$ gives (WB-w).
\item (WB-u) is the same count with Lemma~\ref{lem:wc} (gap $\ge\alpha/4$) in place of Lemma~\ref{lem:sc}, and $n=0$. \qed
\end{itemize}
\end{proof}

\subsection{The mollifier defect on analytic vectors (Lemma UD)}\label{subsec:dd-unw}

\begin{lemma}[UD]\label{lem:unw}
There is a constant $C_0$ such that for $0<\varepsilon\le1$, $0<a\le1$, $s\ge0$ and $z\in\operatorname{Dom}(\ue^{a\Lambda})$,
\[
\norm{(\Delta_\varepsilon(s)-\Delta(s))z}\le C_0\,a^{-4}\,\ue^{-a/(8\varepsilon)}\,\norm{\ue^{a\Lambda}z}.
\]
\end{lemma}

\begin{proof}
\begin{itemize}
\item \emph{The symbol.} For $\lambda\ge0$: $(1+\lambda^2)(1-j_\varepsilon(\lambda))\le5\,a^{-2}\ue^{-a/(4\varepsilon)}\ue^{a\lambda}$. Indeed
$0\le1-j_\varepsilon(\lambda)=\int g(\lambda-\mu)(1-\chi(\varepsilon\mu))\,d\mu\le\int_{|\mu|\ge1/\varepsilon}g(\lambda-\mu)\,d\mu$.
  \begin{itemize}
  \item If $\lambda\le1/(2\varepsilon)$, then $|\lambda-\mu|\ge1/(2\varepsilon)$ on the range of integration, so
  $1-j_\varepsilon(\lambda)\le2\ue^{-1/(8\varepsilon^2)}$. With $u=1/\varepsilon\ge1$ the left side is at most
  $2(1+u^2/4)\ue^{-u^2/8}\le5\ue^{-u/4}\le5\ue^{-a/(4\varepsilon)}$, because $2(1+u^2/4)\ue^{-(u^2-2u)/8}<5$ for $u\ge1$.
  \item If $\lambda>1/(2\varepsilon)$, then $1-j_\varepsilon\le1$ and $(1+\lambda^2)\ue^{-a\lambda/2}\le1+16/(\ue^2a^2)\le4a^{-2}$, so
  $1+\lambda^2\le4a^{-2}\ue^{a\lambda}\ue^{-a\lambda/2}\le4a^{-2}\ue^{a\lambda}\ue^{-a/(4\varepsilon)}$.
  \end{itemize}
\item \emph{The splitting.} $\Delta_\varepsilon(s)-\Delta(s)=-J_\varepsilon\Delta(s)(1-J_\varepsilon)-(1-J_\varepsilon)\Delta(s)$.
\item \emph{First term.} With $z_1:=(1-J_\varepsilon)z$, Lemma~\ref{lem:ell} and $\norm{J_\varepsilon}\le1$ give
$\norm{J_\varepsilon\Delta(s)z_1}\le2(\norm{\Delta z_1}+\norm{z_1})\le4\norm{(1+\Lambda^2)(1-j_\varepsilon(\Lambda))z}\le20\,a^{-2}\ue^{-a/(4\varepsilon)}\norm{\ue^{a\Lambda}z}$,
by the symbol bound.
\item \emph{Second term.} By Lemma~\ref{lem:wc} with $\tau_0=1$ (words of length $1$ and $2$, radii $a/2<a$),
$\norm{\ue^{(a/2)\Lambda}\Delta(s)z}\le3(1+4C_1/a)^2\norm{\ue^{a\Lambda}z}$. By the symbol bound with $a/2$ in place of $a$,
$\norm{(1-J_\varepsilon)w}\le20\,a^{-2}\ue^{-a/(8\varepsilon)}\norm{\ue^{(a/2)\Lambda}w}$. Take $w=\Delta(s)z$.
\end{itemize}
Adding the two terms gives the claim with $C_0=20+60(1+4C_1)^2$. \qed
\end{proof}

\subsection{Regularity and the base point}\label{subsec:dd-base}

The theorems are stated for amplitudes of finite norm and for every boundary point $b$. The next two lemmas reduce them to
$C^\infty$ amplitudes, the base point $i$ and $b=\infty$.

\begin{lemma}[regularity]\label{lem:lift}
Let $0<\alpha\le1$ and $F\in\operatorname{Dom}(\ue^{\alpha\Lambda})$. Then $F$ has a unique $C^\infty$ representative, and every word in
$X,Y$ applied to it lies in $L^2(\Hb)$. Consequently every $a$ with $\norm{\ue^{\alpha\Lambda}X_0^\beta a}<\infty$, and every
$a\in\operatorname{Dom}(\ue^{\alpha\Lambda})$, has a unique $C^\infty$ representative. The approximations $u^{(K)}$ are built from this
representative.
\end{lemma}

\begin{proof}
\begin{itemize}
\item By Lemma~\ref{lem:wc} at $\alpha''=0$ (with $\tau_0=1$), every word applied to $F$ in the sense of distributions lies in $L^2$. The same holds
for $F\circ p$ for every affine isometry $p(z)=p_1z+p_2$, $p_1>0$: these commute with $X$, $Y$ and $\Lambda$, so
$F\circ p\in\operatorname{Dom}(\ue^{\alpha\Lambda})$.
\item Let $w\in L^2$ be such that all words applied to $w$ lie in $L^2$, and let $c\in C_c^\infty(\Hb)$ be real. With the Euclidean
Laplacian $\Delta_{\mathrm E}=\partial_x^2+\partial_y^2=y^{-2}(X^2+Y^2-Y)$,
\[
\Delta_{\mathrm E}(cw)=(\Delta_{\mathrm E}c)\,w+2y^{-2}\bigl((Xc)(Xw)+(Yc)(Yw)\bigr)+c\,y^{-2}(X^2+Y^2-Y)w,
\]
a finite sum of terms $c'w'$ of the same kind. So $\Delta_{\mathrm E}^n(cw)\in L^2(\mathbb R^2)$ for every $n$; on $\operatorname{supp}c$
the measures $dx\,dy$ and $d\nu$ are comparable.
\item Hence $|\xi|^{2n}\widehat{cw}\in L^2(\mathbb R^2)$ for every $n$, all moments of $\widehat{cw}$ are integrable, and $cw$ has a
$C^\infty$ representative.
\item Take $c=1$ near $i$. This gives a $C^\infty$ representative of $F\circ p$ near $i$ for every $p$, that is, of $F$ near
every point. Two continuous representatives on an open set coincide, so they glue to a $C^\infty$ representative of $F$, which is unique.
\item For the consequences take $F=X_0^\beta a$, respectively $F=a$; $X_0^{-\beta}$ is smooth and positive. \qed
\end{itemize}
\end{proof}

\begin{lemma}[base point]\label{lem:base}
Let $x_0\in\Hb$ and $b\in\partial\Hb$.
\begin{itemize}
\item[\textup{(i)}] There is $g\in SL(2,\mathbb R)$, unique up to sign, with $g\,x_0=x_0$ and $g\,\infty=b$. We normalize
$\varphi_b:=\varphi_\infty\circ g^{-1}$. The construction of \cite[(1.20)--(1.29)]{S} is covariant under isometries: the initial state and
the approximations for $b$ are $v\mapsto v\circ g^{-1}$ applied to those for $\infty$ with the amplitude $a\circ g$.
\item[\textup{(ii)}] Let $p(z)=(\operatorname{Im}x_0)\,z+\operatorname{Re}x_0$ and $q:=g\,p$, so $q\,i=x_0$ and $q\,\infty=b$. For an amplitude $a$ put
$a_q:=a\circ q$ and $\Gamma_q:=q^{-1}\Gamma q$. Then, with $c:=\ue^{-i\log(\operatorname{Im}x_0)/\hbar}$,
\[
(a\,\ue^{i\varphi_b/\hbar})\circ q=c\,a_q\,\ue^{i\varphi_\infty/\hbar},\qquad u^{(K)}_b(t)[a]\circ q=c\,u^{(K)}_\infty(t)[a_q],
\]
where the subscript names the boundary point of the construction.
\item[\textup{(iii)}] The map $f\mapsto f\circ q$ is unitary from $L^2(\Hb/\Gamma)$ to $L^2(\Hb/\Gamma_q)$, intertwines the periodizations over $\Gamma$ and
$\Gamma_q$, and commutes with $\Delta$, $\cU(t)$ and $\ue^{\alpha\Lambda}$. The swapped and printed norms of $a$ with base point $x_0$ equal those of $a_q$
with base point $i$. For every $s$, $\sup_z\sum_{\gamma\in\Gamma_q}\ue^{-s\,d(i,\gamma z)}=\sup_z\sum_{\gamma\in\Gamma}\ue^{-s\,d(x_0,\gamma z)}$.
\end{itemize}
Consequently the error norm of $(\Gamma,x_0,b,a)$ equals that of $(\Gamma_q,i,\infty,a_q)$, and an estimate proved for base point $i$ and
$b=\infty$, with constants depending on the group only through quantities invariant under $\Gamma\mapsto\Gamma_q$, $x_0\mapsto i$, holds
for every $x_0$ and $b$.
\end{lemma}

\begin{proof}
\begin{itemize}
\item (i) The stabilizer of $x_0$ in $PSL(2,\mathbb R)$ acts simply transitively on $\partial\Hb$. The phase $\varphi_b$, the transport and
the operators $\Delta(s)$ for $b$ are defined from the geodesics ending at $b$, so they are the transports by $g$ of those for $\infty$.
\item (ii) $\varphi_b\circ q=\varphi_\infty\circ p=\varphi_\infty-\log\operatorname{Im}x_0$, since $\operatorname{Im}p(z)=(\operatorname{Im}x_0)\operatorname{Im}z$.
By (i), composing with $g$ turns the objects for $b$ and $a$ into those for $\infty$ and $a\circ g$. The affine map $p$ commutes with
$X$, $Y$ and $S(t)$, hence with $\Delta(s)$ and the $P_k(t)$, and it changes $\ue^{i\varphi_\infty/\hbar}$ by the constant factor $c$. This
gives both identities.
\item (iii) $q$ is an isometry of $\Hb$ with $q\,i=x_0$, so $d(x_0,qz)=d(i,z)$ and the weights $X_0$, $\langle d\rangle$ based at $x_0$ pull
back to those based at $i$. The spectral functions of $-\Delta$ commute with composition by isometries. For the last identity,
$d(i,q^{-1}\gamma qz)=d(x_0,\gamma(qz))$. \qed
\end{itemize}
\end{proof}

\begin{remark}
The hypothesis $g\,x_0=x_0$ is used in (iii): the Poincar\'e sum of $\Gamma$ is not bounded uniformly in the base point when $\Gamma$
has parabolic elements, so the base point of the weight must be the point at which the sum is taken.
\end{remark}

\subsection{Approximation with the power weight \texorpdfstring{$X_0^\beta$}{X0}}\label{subsec:dd-thm}

\begin{theorem}[approximation with the power weight $X_0^\beta$]\label{thm:dd-appx}
Let $\Gamma$ be any Fuchsian group, $\beta>1/2$ and $\kappa>0$, and $\alpha_0\in(0,\min(1,\pi/4))$. Let $\chi$ be the cutoff of \cite[(4.20)]{S} with the standing hypotheses of Appendix~\ref{app:kl} (real, even, $0\le\chi\le1$).
There are $C,\delta>0$, depending only on $\beta,\kappa,\alpha_0,\chi,\Gamma$, such that the following holds for all $0<\hbar\le1$, all $\alpha\in(0,\alpha_0]$ with $\alpha^3\ge\kappa\hbar$, all $K,N\in\mathbb N_0$, every $b\in\partial\Hb$, every $a$ with $\norm a=\norm{\ue^{\alpha\Lambda}X_0^\beta a}<\infty$, and all $0\le t\le\delta(\alpha/\hbar)^{1/2}$:
\[
\norm{\Delta^N\bigl[\cU(t)(a\,\ue^{i\varphi_b/\hbar})_\Gamma-u^{(K)}(t)_\Gamma\bigr]}_{L^2(M)}\le C^{N+K+1}\,N^{2N}\,K!\,(\alpha\hbar)^{-2N}\Bigl(\frac{\hbar t}{\alpha^2}\Bigr)^{K+1}\norm a,
\]
with $N^{2N}:=1$ at $N=0$.
\end{theorem}

\begin{remark}
This is the power-weight form of the swapped-class estimate. Lemma~\ref{lem:m2w} transfers it to the norm of
Theorem~\ref{thm:dd14}; the deduction follows the proof below. Section~\ref{sec:changes} compares the resulting bounds with \cite{S}.
\end{remark}

\begin{proof}
\medskip\noindent\textbf{Reductions.} By Lemma~\ref{lem:base} we may take $x_0=i$ and $b=\infty$: the group enters the estimates below only
through the Poincar\'e series bound at the base point. By Lemma~\ref{lem:lift}, $a$ is $C^\infty$. We write $\varphi=\varphi_\infty$.

\medskip\noindent\textbf{Parameters.} All lemmas are used with $\tau_0=1$. Let $c_M$ be the constant of Lemma~\ref{lem:mw} and $c_A$ that of
Proposition~\ref{prop:a0}. Put
\[
a_0:=c_M\alpha,\qquad a_1:=\tfrac34a_0,\qquad a_2:=\tfrac12a_0,\qquad\varepsilon:=(\alpha\hbar)^{1/2}/2,\qquad x:=(\alpha/\hbar)^{1/2},
\]
so that $\varepsilon\le\tfrac12$, $\varepsilon/\hbar=x/2$ and $\alpha/\varepsilon=2x$. Fix $\gamma:=(2\beta+1)/4\in(1/2,\beta)$ and put $\theta:=\gamma/\beta\in(0,1)$.

Use (D4) with $\omega=W_\gamma$; every word applied to $a$ lies in $L^2_\omega$ by Lemma~\ref{lem:sc}, since $W_\gamma\le X_0^\beta$. For each $\mathbf s=(s_1,\dots,s_{K+1})\in\Sigma_{K+1}(t)$ let
\[
g(t'):=V_\varepsilon(t',s_1)\,b',\quad0\le t'\le s_1,\qquad w:=g(0)=V_\varepsilon(s_1)^*b',
\]
and write
\[
\mathcal V_{s_1}b'=\mathcal V_0w+\bigl(\mathcal V_{s_1}g(s_1)-\mathcal V_0g(0)\bigr).
\]

\medskip\noindent\textbf{(I) Mollified part.}
\begin{itemize}
\item Lemma~\ref{lem:t2} at $t=0$, and Proposition~\ref{prop:ws} at the weight $X_0^\gamma$ (with $\hbar s_1\le\varepsilon$), give
\[
\norm{\mathcal V_0w}=\norm{[e^{i\varphi/\hbar}V_\varepsilon(s_1)^*b']_\Gamma}\le C\norm{W_\gamma V_\varepsilon(0,s_1)b'}\le C\ue^{G_\gamma/2}\norm{W_\gamma b'}\le C\norm{X_0^\beta b'},
\]
using $W_\gamma=X_0^\gamma\le X_0^\beta$ (since $X_0\ge1$ and $\gamma\le\beta$).
\item By (WB-w) with $n=0$:
\[
\text{(I)}\le C\,C_4^{N+K+1}(2N)!\,(2K+2)!\,\alpha^{-(2K+2)}(\alpha\hbar)^{-2N}\norm a.
\]
\end{itemize}

\medskip\noindent\textbf{(II) Mollifier part.}
\begin{itemize}
\item By (WB-w) with $n\le2$, $X_0^\beta Zb'\in L^2$ for $Z\in\mathcal S=\{1,X,Y,X^2,Y^2\}$. So Lemma~\ref{lem:path}(c), with the exponent
$\gamma$, shows that $g$ satisfies the hypotheses of Lemma~\ref{lem:surf-duh} for $\omega=W_\gamma$. Since $g'=c\,\Delta_\varepsilon(t')g$,
Duhamel on $M$ gives $\mathcal V_{s_1}g(s_1)-\mathcal V_0g(0)=-c\int_0^{s_1}\mathcal V_{t'}d(t')\,dt'$, with
\[
d(t'):=(\Delta(t')-\Delta_\varepsilon(t'))\,g(t').
\]
\item By the unitarity of $\cU(t')$ on $M$ and Lemma~\ref{lem:t2} (exponent $\gamma$),
\[\norm{\mathcal V_{t'}d(t')}=\norm{[e^{i\varphi/\hbar}S(t')d(t')]_\Gamma}\le C\ue^{\gamma t'}\norm{W_\gamma d(t')}.\]
\item H\"older on $L^2$: $\norm{W_\gamma d}\le\norm{W_\beta d}^\theta\norm{d}^{1-\theta}$, since $W_\gamma=W_\beta^\theta$.
\item Weighted factor: by Lemma~\ref{lem:path}(b) with the exponent $\beta$, and (WB-w) with $n\le2$ (five words; the bound increases with $n$),
\begin{align*}
\norm{W_\beta d(t')}&\le C\varepsilon^{-3}\sum_{Z\in\mathcal S}\norm{X_0^\beta Zb'}\\
&\le C\varepsilon^{-3}\,C_4^{N+K+3}(2N)!\,(2K+4)!\,\alpha^{-(2K+4)}(\alpha\hbar)^{-2N}\norm a.
\end{align*}
\item Unweighted factor.
\begin{itemize}
\item Lemma~\ref{lem:mw} gives $\norm{\ue^{a_0\Lambda}a}\le C_*B_\beta\norm a$.
\item (WB-u), applied with $a_0$ in the role of $\alpha$ and the radii $a_1<a_0$ (gap $a_0/4$), gives
$\norm{\ue^{a_1\Lambda}b'}\le C_4^{N+K+1}(2N)!(2K+2)!\,a_0^{-(2K+2)}(a_0\hbar)^{-2N}\norm{\ue^{a_0\Lambda}a}$.
\item Proposition~\ref{prop:a0}, from the radius $a_1$ to $a_2$ (so $(a_1-a_2)/a_1=\tfrac13$) and for $|t'-s_1|\le t\le\tfrac13c_A\varepsilon/\hbar$, gives
$\norm{\ue^{a_2\Lambda}g(t')}\le\norm{\ue^{a_1\Lambda}b'}$.
\item Lemma~\ref{lem:unw} with $a=a_2$ gives $\norm{d(t')}\le C_0a_2^{-4}\ue^{-a_2/(8\varepsilon)}\norm{\ue^{a_2\Lambda}g(t')}$, and $a_2/(8\varepsilon)=c_Mx/8$.
\end{itemize}
\item Combining: put $F:=(C_4/c_M^2)^{N+K+3}(2N)!\,(2K+4)!\,\alpha^{-(2K+4)}(\alpha\hbar)^{-2N}\norm a$. Since $c_M\le1\le C_4$ and $\alpha\le1$, the weighted
factor is at most $C\varepsilon^{-3}F$ and the unweighted one at most $C\alpha^{-4}\ue^{-c_Mx/8}F$. After the H\"older split, using $\varepsilon,\alpha\le1$,
\[
\text{(II)}\le C\,F\,R,\qquad R:=\hbar t\,\ue^{\gamma t}\,\varepsilon^{-3}\,\alpha^{-4}\,\ue^{-(1-\theta)c_Mx/8}.
\]
\item For $t\le\delta x$ with $\delta\le(1-\theta)c_M/(16\gamma)$: $\ue^{\gamma t}\le\ue^{(1-\theta)c_Mx/16}$. Since $\alpha^3\ge\kappa\hbar$ and $\alpha\le1$,
\[
\alpha^{-1}\le x/\sqrt\kappa,\qquad\varepsilon^{-1}=2x/\alpha\le2x^2/\sqrt\kappa,\qquad\hbar t\le\delta\hbar x=\delta\alpha/x\le\delta/x .
\]
Hence $R\le C(\kappa)\,x^{9}\,\ue^{-(1-\theta)c_Mx/16}\le C(\kappa)\,x^{-2}$, and $x^{-2}=\hbar/\alpha\le\alpha^2/\kappa$. So $R\le C(\kappa)\alpha^2$, and
(II) is at most $C\alpha^2F$.
\end{itemize}

\medskip\noindent\textbf{Collecting.}
\begin{itemize}
\item $\int_{\Sigma_{K+1}(t)}ds=t^{K+1}/(K+1)!$.
\item $(2K+4)!/(K+1)!\le C^{K+1}K!$ and $(2K+2)!/(K+1)!\le C^{K+1}K!$, since $(2K+4)!\le4^{K+2}((K+2)!)^2$.
\item $(2N)!\le4^NN^{2N}$.
\item $\hbar^{K+1}t^{K+1}\alpha^{-(2K+2)}=(\hbar t/\alpha^2)^{K+1}$, and $\alpha^2F$ carries $\alpha^{-(2K+2)}$.
\end{itemize}
Hence
\[
\norm E\le C^{N+K+1}N^{2N}K!\,(\alpha\hbar)^{-2N}\Bigl(\frac{\hbar t}{\alpha^2}\Bigr)^{K+1}\norm a.
\]
\qed

\medskip\noindent\textbf{Time windows used.}
\begin{itemize}
\item Weighted propagation and the path bound: $\hbar s_1\le\varepsilon$, that is $t\le x/2$.
\item $A_0$: $t\le\tfrac13c_A\varepsilon/\hbar=c_Ax/6$.
\item Transport: $\ue^{\gamma t}$ is absorbed for $\delta\le(1-\theta)c_M/(16\gamma)$.
\end{itemize}
So $\delta:=\min\bigl(\tfrac12,\;c_A/6,\;(1-\theta)c_M/(16\gamma)\bigr)$ works. The factor $\kappa^{1/2}/\max(1,C)$ is used only in Corollary~\ref{cor:dd15-appx}, to ensure $\delta_K\le1$.

\medskip\noindent\textbf{Dependence on $\Gamma$.} The group enters only through Lemma~\ref{lem:t2}, used in (I) and (II) with exponent $\gamma$, whose
constant $C(\Gamma,\gamma)=2^\gamma P_\Gamma(2\gamma)^{1/2}$ depends on $\Gamma$ only through the Poincar\'e series bound $P_\Gamma(2\gamma)$, $2\gamma=\beta+\tfrac12$.
All other inputs are estimates on $\Hb$. So $\delta$ does not depend on $\Gamma$, and $C$ depends on $\Gamma$ only through $P_\Gamma(\beta+\tfrac12)$;
this is the uniformity claimed in Theorem~\ref{thm:dd14}. (It concerns this theorem only: the extra cap $\delta\le\kappa^{1/2}/\max(1,C)$ of
Corollary~\ref{cor:dd15-appx} depends on $\Gamma$ through $C$.)
\end{proof}

\begin{proof}[Proof of Theorem~\ref{thm:dd14}]
Let $c_0$ and $C'$ be the constants of Lemma~\ref{lem:m2w}, and let $\norm a$ denote the norm of Section~\ref{sec:swapped}, with the weight
$\ue^{\beta\langle d\rangle}$. If $\alpha\le\alpha_0$ and $\alpha^3\ge\kappa\hbar$, then $\bar\alpha:=c_0\alpha\le c_0\alpha_0$,
$\bar\alpha^3\ge c_0^3\kappa\hbar$ and $\norm{\ue^{\bar\alpha\Lambda}X_0^\beta a}\le C'\norm a$. Apply Theorem~\ref{thm:dd-appx} with
$c_0\alpha_0$, $c_0^3\kappa$ and $\bar\alpha$ in place of $\alpha_0$, $\kappa$ and $\alpha$, and let $C_{\rm int},\delta_{\rm int}$ be its constants. Its
window $t\le\delta_{\rm int}(\bar\alpha/\hbar)^{1/2}$ is the window $t\le\delta(\alpha/\hbar)^{1/2}$ with $\delta:=c_0^{1/2}\delta_{\rm int}$, and
$(\bar\alpha\hbar)^{-2N}(\hbar t/\bar\alpha^2)^{K+1}=c_0^{-2N-2K-2}(\alpha\hbar)^{-2N}(\hbar t/\alpha^2)^{K+1}$. So the bound holds with
$C:=C_{\rm int}\,c_0^{-2}\max(1,C')$. The dependence on $\Gamma$ is that of Theorem~\ref{thm:dd-appx}.
\end{proof}

\subsection{Optimal truncation and pointwise estimates}\label{subsec:dd-cor}

\begin{corollary}[optimal truncation with the power weight $X_0^\beta$]\label{cor:dd15-appx}
Under the hypotheses of Theorem~\ref{thm:dd-appx}, with $\delta\le\kappa^{1/2}/\max(1,C)$, let $C$ be its constant, put $\delta_K:=C\hbar t/\alpha^2$, and let $\varepsilon\in(0,1)$. Then for every integer $K$ with $(1-\varepsilon)/\delta_K\le K\le(1-\varepsilon/2)/\delta_K$,
\[
\norm{\Delta^N\bigl[\cU(t)(a\,\ue^{i\varphi_b/\hbar})_\Gamma-u^{(K)}(t)_\Gamma\bigr]}_{L^2(M)}\le C'\varepsilon^{-1/2}C^{N+1}N^{2N}(\alpha\hbar)^{-2N}\ue^{-(1-\varepsilon)\alpha^2/(C\hbar t)}\norm a,
\]
with $C'$ an absolute constant.
\end{corollary}

\begin{proof}[Proof of Corollary~\ref{cor:dd15-appx} (Schubert's computation for his Corollary 1.5 \cite{S})]
\begin{itemize}
\item The bound of Theorem~\ref{thm:dd-appx} is $C^NN^{2N}(\alpha\hbar)^{-2N}\,\delta_K^{K+1}K!\,\norm a$, because $C^{K+1}(\hbar t/\alpha^2)^{K+1}=\delta_K^{K+1}$.
On the window, $K\ge(1-\varepsilon)/\delta_K>0$, so $K\ge1$.
\item \[\delta_K^{K+1}K!\le\delta_K\cdot e\,K^{1/2}\ue^{-K}(K\delta_K)^K\] (Stirling).
\item $K\delta_K\le1-\varepsilon/2$ gives $(K\delta_K)^K\le\ue^{-K\varepsilon/2}$, and $\sup_KK^{1/2}\ue^{-K\varepsilon/2}\le(e\varepsilon)^{-1/2}$.
\item $K\ge(1-\varepsilon)/\delta_K$ gives $\ue^{-K}\le\ue^{-(1-\varepsilon)\alpha^2/(C\hbar t)}$.
\item $\hbar t/\alpha^2\le\delta\,\hbar^{1/2}\alpha^{-3/2}\le\delta\kappa^{-1/2}\le1/\max(1,C)$, by $t\le\delta(\alpha/\hbar)^{1/2}$, $\alpha^3\ge\kappa\hbar$ and
$\delta\le\kappa^{1/2}/\max(1,C)$. Hence $\delta_K\le1$ and $C^N\delta_K=C^{N+1}\hbar t/\alpha^2\le C^{N+1}$, and the bound holds with $C'=\ue^{1/2}$. \qed
\end{itemize}
\end{proof}

\begin{remark}
In \cite[Corollary~1.5]{S} the window is printed as $K\in[C\alpha^2/\hbar t-\varepsilon,\;C\alpha^2/\hbar t-\varepsilon/2]$, with the constant $C$ of
\cite[(1.30)]{S}. The proof there uses $K\delta_K\in[1-\varepsilon,1-\varepsilon/2]$, which is the window used here. If the window contains no integer
the statement is empty; at $t=0$ the error is $0$.
\end{remark}

\begin{corollary}[pointwise approximation with the power weight $X_0^\beta$]\label{cor:dd16-appx}
Assume, in addition, that $M$ has injectivity radius $\ge r_0>0$. Then, under the hypotheses of Theorem~\ref{thm:dd-appx}, the error
$\cU(t)(a\,\ue^{i\varphi_b/\hbar})_\Gamma-u^{(K)}(t)_\Gamma$ lies in $\operatorname{Dom}(\Delta_M^N)$, its image under $\Delta^N$ has a continuous
representative, and at every point of $M$
\[
\Bigl|\Delta^N\bigl[\cU(t)(a\,\ue^{i\varphi_b/\hbar})_\Gamma-u^{(K)}(t)_\Gamma\bigr]\Bigr|\le C^{N+K+1}(2N)!\,K!\,(\alpha\hbar)^{-2N-4}\Bigl(\frac{\hbar t}{\alpha^2}\Bigr)^{K+1}\norm a,
\]
with $C$ depending also on $r_0$. In the same way Corollary~\ref{cor:dd15-appx} gives the pointwise exponential bound with $(\alpha\hbar)^{-2N-4}$.
\end{corollary}

\begin{proof}
\begin{itemize}
\item If $d(\tilde z,\gamma\tilde z)\ge2r_0$ for all $\gamma\ne1$, the balls $B(\gamma\tilde z,r_0)$ are disjoint, and those with
$d(\tilde z,\gamma\tilde z)<2$ lie in $B(\tilde z,2+r_0)$. So the multiplicity of Lemma~\ref{pc:g43} satisfies $m(z)\le C(r_0)$, and that lemma gives, for
$u\in\operatorname{Dom}(\Delta_M)$, a continuous representative with
$\sup_M|u|\le C(r_0)\bigl(\norm{\Delta u}_{L^2(M)}+\norm u_{L^2(M)}\bigr)$.
\item Apply this to $u=\Delta^NE$, with Theorem~\ref{thm:dd-appx} at $N$ and at $N+1$. Use $N^{2N}\le(\ue/2)^{2N}(2N)!$,
$(N+1)^{2N+2}\le54^{N+1}(2N)!$ and $\alpha\hbar\le1$. \qed
\end{itemize}
\end{proof}

\begin{remark}
The proof uses one Laplacian and gives the exponent $-2N-2$ of $\alpha\hbar$. We state the corollary with $-2N-4$, the exponent that
enters the comparison with \cite[(1.34)]{S} in Section~\ref{sec:changes}.
\end{remark}

\begin{remark}[Cusped $M$]
The Sobolev inequality \cite[(5.91)]{S} fails on $M$ with cusps (Remark~\ref{rem:g4}). On such $M$ the same proof, with Lemma~\ref{pc:g43} as it stands, gives the pointwise bound with the additional factor $C\,m(z)^{1/2}$, as in Corollary~\ref{pc:cor16}.
\end{remark}

% The corollary body ends before the remarks.

\subsection{Lemma EX: the example amplitudes of \texorpdfstring{\cite{S}}{[S]} are covered}\label{subsec:dd-ex}

\begin{lemma}[EX]\label{lem:ex}
Fix $r_1<\tfrac12$. There are constants $C_H$ and $C_{\mathrm E}$, depending only on $r_1$, such that for $\mu\ge1$ and
$0<\alpha\le\min(r_1,\mu^{-1/2})$,
\[
\norm{\ue^{\alpha\Lambda}\ue^{-\mu\langle d\rangle}}_{L^2(\Hb)}\le(\pi\ue)^{1/2}\,\ue^{C_H\mu\alpha^2}\norm{\ue^{-\mu\langle d\rangle}}_{L^2(\Hb)}\le C_{\mathrm E}\,\mu^{-1/2}.
\]
\end{lemma}

\begin{proof}
We use the notation of Section~\ref{subsec:gutzmer}, with $x_0=i$.

\emph{The function $H$.} Let $G(z):=-\arccos(z)^2$ and $H(z):=(1+G(z))^{1/2}-1$. The function $G$ is holomorphic on
$\mathbb C\setminus(-\infty,-1]$ and equals $d^2$ at $z=\cosh d$, so $\langle d\rangle=H(X_0)$ on $\Hb$. The number $1+G(z)$ is real and at most $0$ only for
real $z\le\cos1$, so $H$ is holomorphic on the complement of $(-\infty,\cos1]$, real on $(\cos1,\infty)$ and nondecreasing there.
For $0<\rho<\tfrac12$ let $S_\rho:=\{\operatorname{Re}z\ge\cos2\rho,\ |\arg z|\le2\rho\}$, and fix $r_1<r_1'<r_2<\tfrac12$ depending only on $r_1$.
Since $\cos2r_2>\cos1$, $H$ is holomorphic near $S_{r_2}$. As $|z|\to\infty$ in $S_{r_2}$ we have $G=L^2$ with
$L=\log(z+(z^2-1)^{1/2})=\log(2z)+O(|z|^{-2})$, hence $H=L-1+O(1/|L|)$; so $M_H:=\sup_{S_{r_2}}|H(z)-\log z|<\infty$. There is $c>0$,
depending only on $r_1',r_2$, such that the disc of radius $c|z|$ about any $z\in S_{r_1'}$ lies in $S_{r_2}$, and Cauchy's estimates for
$H-\log$ on these discs give
\[
|H'(z)|\le C_1/|z|,\qquad|H''(z)|\le C_1/|z|^2\qquad(z\in S_{r_1'}).
\]

\emph{The orbit norms.} By Lemma~\ref{lem:gutzmer}(b), on the orbit $G\exp(i\rho h_0)x_0$, $\rho\le r_1'$, the extension of $X_0$ takes the
values $\tilde X_0=cA+isB$ with $c=\cos2\rho$, $s=\sin2\rho$, $A=X_0(gx_0)\ge1$ and $|B|\le A$; these lie in $S_\rho\subset S_{r_1'}$, and so does
the vertical segment from $cA$ to $\tilde X_0$. So $f:=\ue^{-\mu\langle d\rangle}=\ue^{-\mu H(X_0)}$ extends holomorphically to $\Xi_{r_1'}$. Since $H'(cA)$ is
real, Taylor's formula on that segment gives $\operatorname{Re}H(\tilde X_0)\ge H(cA)-\tfrac12\sup|H''|\,s^2B^2\ge H(cA)-C_1s^2/(2c^2)$, and
$H(A)-H(cA)\le(1-c)A\sup_{[cA,A]}|H'|\le C_1(1-c)/c$. With $1-c\le2\rho^2$, $s^2\le4\rho^2$ and $c\ge\cos1>\tfrac12$,
\[
\operatorname{Re}H(\tilde X_0)\ge\langle d\rangle(gx_0)-C_H\rho^2,\qquad C_H:=12\,C_1 .
\]
Hence $|f(g\exp(i\rho h_0)x_0)|\le\ue^{C_H\mu\rho^2}f(gx_0)$ and $N_\rho(f)\le\ue^{C_H\mu\rho^2}\norm f$ for $\rho\le r_1'$. In particular
$\sup_{\rho<r_1'}N_\rho(f)<\infty$, and Lemma~\ref{lem:gutzmer}(a), with $R=r_1'$ and $r=\alpha\le r_1$, gives the first inequality.

\emph{The $L^2$ norm.} In geodesic polar coordinates $\norm f^2=2\pi\int_0^\infty\ue^{-2\mu\langle\rho\rangle}\sinh\rho\,d\rho$, with
$\langle\rho\rangle=(1+\rho^2)^{1/2}-1$. For $\rho\le1$, $\langle\rho\rangle\ge\rho^2/(1+\sqrt2)$ and $\sinh\rho\le2\rho$; for $\rho\ge1$,
$\langle\rho\rangle\ge\rho-1$ and $\sinh\rho\le\ue^\rho$. So for $\mu\ge1$
\[
\norm f^2\le2\pi\Bigl(\int_0^\infty2\rho\,\ue^{-2\mu\rho^2/(1+\sqrt2)}d\rho+\ue\int_0^\infty\ue^{-(2\mu-1)u}du\Bigr)\le\frac{2\pi(2+\ue)}{\mu},
\]
and the second inequality holds with $C_{\mathrm E}:=(\pi\ue)^{1/2}\ue^{C_H}(2\pi(2+\ue))^{1/2}$, because $\mu\alpha^2\le1$.
\end{proof}

\begin{remark}[Consequence for the example of \texorpdfstring{\cite{S}}{[S]}]
\cite[Remark~1.7, (1.38)]{S} gives the example $a_\hbar:=\hbar^{-\delta/2}\ue^{-\hbar^{-\delta}\langle d\rangle}$, said there to lie in $H_{\alpha,\beta}$ with $\beta=1$ and $\alpha=\hbar^\delta$, and to give long times for $\delta<1/3$. In the swapped $\langle d\rangle$-class, $\norm{a_\hbar}^{\rm sw}_{\alpha,\beta}=\hbar^{-\delta/2}\norm{\ue^{\alpha\Lambda}\ue^{-(\hbar^{-\delta}-\beta)\langle d\rangle}}$. Apply Lemma~\ref{lem:ex} with $\mu=\hbar^{-\delta}-\beta$ and $\alpha_\hbar:=(\hbar^{-\delta}-\beta)^{-1/2}\asymp\hbar^{\delta/2}$:
\begin{itemize}
\item $\norm{a_\hbar}^{\rm sw}_{\alpha_\hbar,\beta}\le C$ uniformly in $\hbar$, for every fixed $\beta$ (not only $\beta=1$);
\item the swapped $\langle d\rangle$-class transfers to the $\tilde\psi$-class by Lemma~\ref{lem:m2w}, at the radius $c_0\alpha_\hbar$ with the fixed ratio $c_0$;
\item so Theorem~\ref{thm:dd-appx} (any $\beta>1/2$, every Fuchsian $\Gamma$) applies when $(c_0\alpha_\hbar)^3\ge\kappa\hbar$, i.e. for $0<\delta<2/3$, and for $\hbar$ small enough that $\mu=\hbar^{-\delta}-\beta\ge1$ and $\alpha_\hbar\le r_1$. The case $\delta=0$ is not covered: $\ue^{-\langle d\rangle}$ has infinite $\ue^{\beta\langle d\rangle}$-weighted $L^2$ norm for every $\beta>1/2$. The time window is $t\le\delta'(\alpha_\hbar/\hbar)^{1/2}\asymp\hbar^{\delta/4-1/2}$.
\end{itemize}
The radius $\alpha=\hbar^\delta$ of \cite{S} gave $\delta<1/3$ and $t\ll\hbar^{-(1-3\delta)/2}$; the holomorphic tube bound gives the larger radius $\hbar^{\delta/2}$. The example is therefore covered. The printed (original-order) class is treated in Section~\ref{sec:printed}; the amplitudes named in \cite{S} are covered already by the swapped class.
\end{remark}

\section{Periodization by Cauchy--Schwarz}\label{sec:t2}

\begin{lemma}[P1]\label{lem:p1}
For every Fuchsian $\Gamma$ and $s>1$: $\sup_{z}\sum_{\gamma}\ue^{-s\,d(x_0,\gamma z)}<\infty$.
\end{lemma}
\begin{proof}
The orbit $\Gamma x_0$ is uniformly separated, and its stabilizer is finite. Packing balls of a fixed radius
around orbit points gives $\#\{\gamma:d(z,\gamma x_0)<n+1\}\le C\ue^{n}$, uniformly in $z$.
\end{proof}

\begin{lemma}[T2]\label{lem:t2}
For every Fuchsian $\Gamma$, every $s>1/2$, $t\ge0$, every real phase $\varphi$ and every measurable $w$ with $\ue^{s\tilde\psi}w\in L^2(\Hb)$:
\[
\norm{\bigl[(S(t)w)\,\ue^{i\varphi/\hbar}\bigr]_\Gamma}_{L^2(M)}\le C(\Gamma,s)\,\ue^{st}\,\norm{\ue^{s\tilde\psi}w}_{L^2(\Hb)},
\]
with $C(\Gamma,s)=2^sP_\Gamma(2s)^{1/2}$, where $P_\Gamma(\sigma):=\sup_z\sum_{\gamma\in\Gamma}\ue^{-\sigma d(x_0,\gamma z)}$ is the Poincar\'e series bound of Lemma~\ref{lem:p1}.
\end{lemma}

\begin{proof}
\begin{itemize}
\item \emph{Pointwise Cauchy--Schwarz.} For $f$ on $\Hb$ and $z\in\Hb$:
\[
\Bigl|\sum_\gamma f(\gamma z)\Bigr|^2\le\Bigl(\sum_\gamma\ue^{-2sd(\gamma z,x_0)}\Bigr)\Bigl(\sum_\gamma\ue^{2sd(\gamma z,x_0)}|f(\gamma z)|^2\Bigr).
\]
The first factor is $\le P_\Gamma(2s)$ for every $z$, by Lemma~\ref{lem:p1} with exponent $2s>1$ (since $d(\gamma z,x_0)=d(z,\gamma^{-1}x_0)$).
\item \emph{Unfolding.} Integrate over a fundamental domain (Lemma 3.1 of \cite{S}, Tonelli). This gives
\[
\norm{[f]_\Gamma}^2_{L^2(M)}\le P_\Gamma(2s)\int_\Hb\ue^{2sd(z,x_0)}|f(z)|^2\,d\nu(z).
\]
Prove this first for finite partial sums, then pass to the limit by monotone convergence. In particular $\sum_\gamma|f\circ\gamma|$ converges a.e. and in $L^2(M)$. Torsion-freeness is not needed: the fixed sets of elliptic elements are null, so unfolding over a fundamental domain carries no stabilizer factor, and any multiplicity is already counted in the orbit series $P_\Gamma(2s)$.
\item \emph{Transport.} Take $f=(S(t)w)\ue^{i\varphi/\hbar}$, so $|f|=|S(t)w|$.
  \begin{itemize}
  \item $(S(t)w)(z)=\ue^{-t/2}w(z_t)$, where $z_t$ is at distance $t$ from $z$, so $d(z,x_0)\le d(z_t,x_0)+t$.
  \item $S(t)$ is unitary, so $\int\ue^{2sd(z)}|S(t)w(z)|^2\,d\nu(z)\le\ue^{2st}\int\ue^{2sd(z')}|w(z')|^2\,d\nu(z')$.
  \end{itemize}
\item Finally $\ue^{sd}\le2^s\ue^{s\tilde\psi}$, since $\tilde\psi=\log\cosh d\ge d-\log2$. \qed
\end{itemize}
\end{proof}

% End of dd_appendix.tex
